% Generated from lessons/0012-4-implicit-functions.html; edit the HTML source, then rerun the converter.
\section{隐函数}
\lessonmark{隐函数}

这一节回答一个很现实的问题：方程 \(F(x,y)=0\) 没有显式写出 \(y=f(x)\)，但在某个点附近能不能把 \(y\) 看成 \(x\) 的函数？如果能，它的导数怎么求？

\subsection*{本节主线}

\begin{conceptbox}
\textbf{存在唯一}
方程 \(F(x,y)=0\) 附近能不能唯一解出 \(y=f(x)\)。
\end{conceptbox}

\begin{conceptbox}
\textbf{求导公式}
不用真的解出 \(f\)，直接对方程两边求导。
\end{conceptbox}

\begin{conceptbox}
\textbf{Jacobi 判别}
多个方程时，看关于未知函数的 Jacobi 行列式是否非零。
\end{conceptbox}

\begin{notebox}
一句话：隐函数题先问“能不能局部解出”，再问“解出来后怎么求导”。考试计算通常只需要后半句：对方程求导，解线性方程组。
\end{notebox}

\subsection*{一、单个方程：从圆的例子看隐函数}

单位圆

\[
      x^2+y^2=1
    \]

整体上不能把 \(y\) 看成 \(x\) 的单值函数，因为同一个 \(x\) 可能对应上下两个 \(y\)。但在上半圆附近可以写成

\[
      y=\sqrt{1-x^2}.
    \]

在下半圆附近可以写成

\[
      y=-\sqrt{1-x^2}.
    \]

所以隐函数定理是一个局部定理：它不保证全局能解，只保证在某个点附近能唯一解出。

\subsection*{二、定理 12.4.1：一元隐函数存在定理}

设 \(F(x,y)\) 满足：

\begin{enumerate}
 \item \(F(x_0,y_0)=0\)；
 \item \(F\) 在矩形邻域内连续，并且 \(\frac{\partial F}{\partial x},\frac{\partial F}{\partial y}\) 连续；
 \item \(\frac{\partial F}{\partial y}(x_0,y_0)\ne0\)。
 \end{enumerate}

那么在 \(x_0\) 附近，可以由方程

\[
      F(x,y)=0
    \]

唯一确定一个隐函数

\[
      y=f(x),\qquad f(x_0)=y_0,
    \]

并且 \(f\) 连续、可导，导数为

\[
      \frac{dy}{dx}
      =
      f'(x)
      =
      -\frac{\frac{\partial F}{\partial x}(x,f(x))}
      {\frac{\partial F}{\partial y}(x,f(x))}.
    \]

\subsubsection*{证明思路}

若 \(\frac{\partial F}{\partial y}(x_0,y_0)>0\)，则在小矩形里 \(\frac{\partial F}{\partial y}>0\)，所以固定 \(x\) 时，\(F(x,y)\) 关于 \(y\) 严格增加。再用介值定理保证有解，用严格单调保证唯一。

可导性来自对恒等式 \(F(x,f(x))=0\) 求导：

\[
      \frac{d}{dx}F(x,f(x))
      =
      \frac{\partial F}{\partial x}
      +
      \frac{\partial F}{\partial y}\frac{dy}{dx}
      =
      0.
    \]

\subsection*{三、定理 12.4.2：多元自变量，一个未知函数}

若

\[
      F(x_1,\ldots,x_n,y)=0
    \]

在点 \((x_1^0,\ldots,x_n^0,y^0)\) 附近满足

\[
      F(x_1^0,\ldots,x_n^0,y^0)=0,\qquad
      \frac{\partial F}{\partial y}(x_1^0,\ldots,x_n^0,y^0)\ne0,
    \]

且相关偏导连续，则局部唯一确定

\[
      y=f(x_1,\ldots,x_n).
    \]

求偏导时仍然对方程两边求导：

\[
      \frac{\partial y}{\partial x_i}
      =
      -\frac{\frac{\partial F}{\partial x_i}}
      {\frac{\partial F}{\partial y}},
      \qquad i=1,\ldots,n.
    \]

\begin{notebox}
这条公式的记法：分母永远是“被解出的那个变量”的偏导；分子是你正在对哪个自变量求导。
\end{notebox}

\subsection*{四、多个方程：用 Jacobi 行列式判断}

\subsubsection*{定理 12.4.3：两个方程、两个未知函数}

设方程组

\[
      \begin{cases}
      F(x,y,u,v)=0,\\
      G(x,y,u,v)=0
      \end{cases}
    \]

希望在某点附近把 \(u,v\) 看成 \(x,y\) 的函数：

\[
      u=f(x,y),\qquad v=g(x,y).
    \]

关键条件是关于未知量 \(u,v\) 的 Jacobi 行列式非零：

\[
      \frac{\partial(F,G)}{\partial(u,v)}
      =
      \begin{vmatrix}
      \frac{\partial F}{\partial u} & \frac{\partial F}{\partial v}\\
      \frac{\partial G}{\partial u} & \frac{\partial G}{\partial v}
      \end{vmatrix}
      \ne0.
    \]

导数矩阵满足

\[
      \begin{pmatrix}
      \frac{\partial u}{\partial x} & \frac{\partial u}{\partial y}\\
      \frac{\partial v}{\partial x} & \frac{\partial v}{\partial y}
      \end{pmatrix}
      =
      -
      \begin{pmatrix}
      \frac{\partial F}{\partial u} & \frac{\partial F}{\partial v}\\
      \frac{\partial G}{\partial u} & \frac{\partial G}{\partial v}
      \end{pmatrix}^{-1}
      \begin{pmatrix}
      \frac{\partial F}{\partial x} & \frac{\partial F}{\partial y}\\
      \frac{\partial G}{\partial x} & \frac{\partial G}{\partial y}
      \end{pmatrix}.
    \]

\subsubsection*{定理 12.4.4：一般 \(m\) 个方程}

对方程组

\[
      F_i(x_1,\ldots,x_n,y_1,\ldots,y_m)=0,\qquad i=1,\ldots,m,
    \]

若关于未知量 \(y_1,\ldots,y_m\) 的 Jacobi 行列式非零：

\[
      \frac{\partial(F_1,\ldots,F_m)}{\partial(y_1,\ldots,y_m)}\ne0,
    \]

则局部唯一确定

\[
      y_j=f_j(x_1,\ldots,x_n),\qquad j=1,\ldots,m.
    \]

\subsection*{五、逆映射定理}

若映射

\[
      x=x(u,v),\qquad y=y(u,v)
    \]

在点 \((u_0,v_0)\) 附近有连续偏导，并且 Jacobi 行列式

\[
      \frac{\partial(x,y)}{\partial(u,v)}\ne0,
    \]

则局部存在反函数

\[
      u=u(x,y),\qquad v=v(x,y).
    \]

反函数的 Jacobi 矩阵是原 Jacobi 矩阵的逆矩阵。

\[
      \frac{\partial(u,v)}{\partial(x,y)}
      =
      \left(
      \frac{\partial(x,y)}{\partial(u,v)}
      \right)^{-1}.
    \]

特别地，在二维行列式层面：

\[
      \frac{\partial(u,v)}{\partial(x,y)}
      =
      \frac{1}{\frac{\partial(x,y)}{\partial(u,v)}}.
    \]

\begin{notebox}
极坐标变换中 \(\frac{\partial(x,y)}{\partial(r,\theta)}=r\)。当 \(r\ne0\) 时，局部可以反解出 \(r,\theta\)。原点处 \(r=0\)，角度不唯一，所以不能局部反解。
\end{notebox}

\subsection*{六、考试做题流程}

\begin{center}
\small
\begin{tabularx}{\linewidth}{|>{\raggedright\arraybackslash}p{0.18\linewidth}|>{\raggedright\arraybackslash}p{0.42\linewidth}|>{\raggedright\arraybackslash}X|}
\hline
\textbf{题型} & \textbf{动作} & \textbf{注意点} \\ \hline
单方程 \(F(x,y)=0\) & 写 \(\frac{\partial F}{\partial x}+\frac{\partial F}{\partial y}\frac{dy}{dx}=0\)。 & 分母是 \(\frac{\partial F}{\partial y}\)，不能为零。 \\ \hline
单方程 \(F(x,y,z)=0\) & 若解 \(z=z(x,y)\)，分别对 \(x,y\) 求偏导。 & \(\frac{\partial z}{\partial x}=-F_x/F_z,\ \frac{\partial z}{\partial y}=-F_y/F_z\)。 \\ \hline
方程组 & 把未知函数的导数组成线性方程组。 & 用关于未知量的 Jacobi 行列式判断可解。 \\ \hline
二阶偏导 & 先求一阶，再对一阶关系继续求导。 & 不要忘记 \(z=z(x,y)\) 或 \(y=y(x)\) 仍然是函数。 \\ \hline
\end{tabularx}
\end{center}

\subsection*{七、即时判断练习}

\begin{quickcheck}
若 \(F(x,y)=0\) 要局部解出 \(y=f(x)\)，关键非零条件是哪一个？

\tcblower
\textbf{答案：}b
\end{quickcheck}

\begin{quickcheck}
方程组想解出 \(u,v\) 为 \(x,y\) 的函数，应检查哪个 Jacobi 行列式？

\tcblower
\textbf{答案：}a
\end{quickcheck}

\subsection*{八、课后题训练}

下面按教材第 165-167 页习题 1-14 整理。题面完整保留，提示只给入口；写作业时仍建议先独立算一遍。

\begin{exercisebox}
\textbf{习题 1\badge{单方程隐函数求导}}

求下列方程所确定的隐函数的导数或偏导数：

\[
      \begin{aligned}
      &(1)\ \sin y+e^x-xy^2=0,\quad \text{求}\ \frac{dy}{dx};\\
      &(2)\ x^y=y^x,\quad \text{求}\ \frac{dy}{dx};\\
      &(3)\ \ln\sqrt{x^2+y^2}=\arctan\frac{y}{x},\quad \text{求}\ \frac{dy}{dx};\\
      &(4)\ \arctan\frac{x+y}{a}-\frac{y}{a}=0,\quad \text{求}\ \frac{dy}{dx}\ \text{和}\ \frac{d^2y}{dx^2};\\
      &(5)\ \frac{x}{z}=\ln\frac{z}{y},\quad \text{求}\ \frac{\partial z}{\partial x}\ \text{和}\ \frac{\partial z}{\partial y};\\
      &(6)\ e^z-xyz=0,\quad \text{求}\ \frac{\partial z}{\partial x},\frac{\partial z}{\partial y},
      \frac{\partial^2 z}{\partial x^2},\frac{\partial^2 z}{\partial x\partial y};\\
      &(7)\ z^3-3xyz=a^3,\quad \text{求}\ \frac{\partial z}{\partial x},\frac{\partial z}{\partial y},
      \frac{\partial^2 z}{\partial x^2},\frac{\partial^2 z}{\partial x\partial y};\\
      &(8)\ f(x+y,y+z,z+x)=0,\quad \text{求}\ \frac{\partial z}{\partial x}\ \text{和}\ \frac{\partial z}{\partial y};\\
      &(9)\ z=f(xz,z-y),\quad \text{求}\ \frac{\partial z}{\partial x},\frac{\partial z}{\partial y},\frac{\partial^2 z}{\partial x^2};\\
      &(10)\ f(x,x+y,x+y+z)=0,\quad \text{求}\ \frac{\partial z}{\partial x},\frac{\partial z}{\partial y},
      \frac{\partial^2 z}{\partial x^2},\frac{\partial^2 z}{\partial x\partial y}.
      \end{aligned}
      \]

\begin{hintbox}
统一写成 \(F=0\)。一阶偏导直接用 \(\frac{\partial z}{\partial x}=-F_x/F_z\)；二阶偏导不要跳步，先保留一阶关系再继续求偏导。
\end{hintbox}
\end{exercisebox}

\begin{exercisebox}
\textbf{习题 2\badge{高阶导验证}}

设 \(y=\tan(x+y)\) 确定 \(y\) 为 \(x\) 的隐函数，验证

\[
        \frac{d^3y}{dx^3}
        =
        \frac{2(3y^4+8y^2+5)}{y^7}.
      \]

\begin{hintbox}
先从 \(\tan(x+y)=y\) 得到一阶导，再连续求导。每一步都把 \(\sec^2(x+y)\) 用 \(1+\tan^2(x+y)=1+y^2\) 消去。
\end{hintbox}
\end{exercisebox}

\begin{exercisebox}
\textbf{习题 3\badge{构造型隐函数}}

设 \(\varphi\) 是可微函数，证明由

\[
        \varphi(cx-az,cy-bz)=0
      \]

所确定的隐函数 \(z=f(x,y)\) 满足

\[
        a\frac{\partial z}{\partial x}
        +
        b\frac{\partial z}{\partial y}
        =
        c.
      \]

\begin{hintbox}
分别对 \(x,y\) 求偏导，得到两个含 \(\varphi_1,\varphi_2\) 的方程，再线性组合消去它们。
\end{hintbox}
\end{exercisebox}

\begin{exercisebox}
\textbf{习题 4\badge{齐次结构}}

设方程

\[
        \varphi\left(x+\frac{z}{y},y+\frac{z}{x}\right)=0
      \]

确定隐函数 \(z=f(x,y)\)，证明它满足

\[
        x\frac{\partial z}{\partial x}
        +
        y\frac{\partial z}{\partial y}
        =
        z-xy.
      \]

\begin{hintbox}
对 \(x,y\) 分别求偏导后，尝试做 \(x\) 倍第一式加 \(y\) 倍第二式。
\end{hintbox}
\end{exercisebox}

\begin{exercisebox}
\textbf{习题 5\badge{方程组隐函数}}

求下列方程组所确定的隐函数的导数或偏导数：

\[
      \begin{aligned}
      &(1)\ \begin{cases}
      z-x^2-y^2=0,\\
      x^2+2y^2+3z^2=4a^2,
      \end{cases}
      \quad \text{求}\ \frac{dy}{dx},\frac{dz}{dx},\frac{d^2y}{dx^2},\frac{d^2z}{dx^2};\\
      &(2)\ \begin{cases}
      xu+yv=0,\\
      yu+xv=1,
      \end{cases}
      \quad \text{求}\ \frac{\partial u}{\partial x},\frac{\partial u}{\partial y},
      \frac{\partial^2u}{\partial x^2},\frac{\partial^2u}{\partial x\partial y};\\
      &(3)\ \begin{cases}
      u=f(ux,v+y),\\
      v=g(u-x,v^2y),
      \end{cases}
      \quad \text{求}\ \frac{\partial u}{\partial x}\ \text{和}\ \frac{\partial v}{\partial x};\\
      &(4)\ \begin{cases}
      x=u+v,\\
      y=u-v,\\
      z=u^2v^2,
      \end{cases}
      \quad \text{求}\ \frac{\partial z}{\partial x}\ \text{和}\ \frac{\partial z}{\partial y};\\
      &(5)\ \begin{cases}
      x=e^u\cos v,\\
      y=e^u\sin v,\\
      z=u^2+v^2,
      \end{cases}
      \quad \text{求}\ \frac{\partial z}{\partial x}\ \text{和}\ \frac{\partial z}{\partial y}.
      \end{aligned}
      \]

\begin{hintbox}
先分清哪些是自变量、哪些是未知函数。方程组求导后得到线性方程组，用 Jacobi 行列式或矩阵法解。
\end{hintbox}
\end{exercisebox}

\begin{exercisebox}
\textbf{习题 6\badge{求微分}}

求微分：

\[
      \begin{aligned}
      &(1)\ x+2y+z-2\sqrt{xyz}=0,\quad \text{求}\ dz;\\
      &(2)\ \begin{cases}
      x+y=u+v,\\
      \frac{x}{\sin u}=\frac{y}{\sin v},
      \end{cases}
      \quad \text{求}\ du\ \text{与}\ dv.
      \end{aligned}
      \]

\begin{hintbox}
把 \(dx,dy,dz,du,dv\) 当作一阶线性关系中的未知量或已知量处理。
\end{hintbox}
\end{exercisebox}

\begin{exercisebox}
\textbf{习题 7\badge{向量值隐函数}}

设

\[
        x=x(y),\qquad z=z(y)
      \]

是由方程组

\[
        F(x,y,z)=0,\qquad
        G\left(\frac{x}{y},\frac{z}{y}\right)=0
      \]

所确定的向量值隐函数，其中二元函数 \(F\) 和 \(G\) 均具有连续偏导数，求

\[
        \frac{dx}{dy}\quad \text{和}\quad \frac{dz}{dy}.
      \]

\begin{hintbox}
对 \(y\) 求导，得到关于 \(dx/dy,dz/dy\) 的二元线性方程组。
\end{hintbox}
\end{exercisebox}

\begin{exercisebox}
\textbf{习题 8\badge{极坐标 Laplace 算子}}

设 \(f(x,y)\) 具有二阶连续偏导数。在极坐标

\[
        x=r\cos\theta,\qquad y=r\sin\theta
      \]

变换下，求

\[
        \frac{\partial^2 f}{\partial x^2}
        +
        \frac{\partial^2 f}{\partial y^2}
      \]

关于极坐标的表达式。

\begin{hintbox}
目标结果是二维 Laplace 算子的极坐标形式：径向二阶项、径向一阶项和角向二阶项。
\end{hintbox}
\end{exercisebox}

\begin{exercisebox}
\textbf{习题 9\badge{线性变换化简 PDE}}

设二元函数 \(f\) 具有二阶连续偏导数。证明：通过适当线性变换

\[
        u=x+\lambda y,\qquad v=x+\mu y,
      \]

可以将方程

\[
        A\frac{\partial^2 f}{\partial x^2}
        +
        2B\frac{\partial^2 f}{\partial x\partial y}
        +
        C\frac{\partial^2 f}{\partial y^2}=0,\qquad AC-B^2<0
      \]

化简为

\[
        \frac{\partial^2 f}{\partial u\partial v}=0.
      \]

并说明此时 \(\lambda,\mu\) 为二元二次方程 \(A+2Bt+Ct^2=0\) 的两个相异实根。
\end{exercisebox}

\begin{exercisebox}
\textbf{习题 10\badge{指数变量替换}}

通过自变量变换

\[
        x=e^\xi,\qquad y=e^\eta,
      \]

变换方程

\[
        ax^2\frac{\partial^2 z}{\partial x^2}
        +
        2bxy\frac{\partial^2 z}{\partial x\partial y}
        +
        cy^2\frac{\partial^2 z}{\partial y^2}=0,
      \]

其中 \(a,b,c\) 为常数。
\end{exercisebox}

\begin{exercisebox}
\textbf{习题 11\badge{开方变量替换}}

通过自变量变换

\[
        u=x-2\sqrt y,\qquad v=x+2\sqrt y,
      \]

变换方程

\[
        \frac{\partial^2 z}{\partial x^2}
        -
        y\frac{\partial^2 z}{\partial y^2}
        =
        \frac12\frac{\partial z}{\partial y},\qquad y>0.
      \]
\end{exercisebox}

\begin{exercisebox}
\textbf{习题 12\badge{同时换自变量和因变量}}

导出新的因变量关于新的自变量的偏导数所满足的方程：

\[
      \begin{aligned}
      &(1)\ \text{用}
      \begin{cases}
      u=x^2+y^2,\\
      v=\frac1x+\frac1y,
      \end{cases}
      \quad \text{及}\quad w=\ln z-(x+y)
      \quad \text{变换方程}\\
      &\qquad
      y\frac{\partial z}{\partial x}
      -
      x\frac{\partial z}{\partial y}
      =
      (y-x)z;\\[4pt]
      &(2)\ \text{用}
      \begin{cases}
      u=x,\\
      v=x+y,
      \end{cases}
      \quad \text{及}\quad w=x+y+z
      \quad \text{变换方程}\\
      &\qquad
      \frac{\partial^2 z}{\partial x^2}
      -
      2\frac{\partial^2 z}{\partial x\partial y}
      +
      \left(1+\frac{y}{x}\right)
      \frac{\partial^2 z}{\partial y^2}=0;\\[4pt]
      &(3)\ \text{用}
      \begin{cases}
      u=x+y,\\
      v=\frac{y}{x},
      \end{cases}
      \quad \text{及}\quad w=\frac{z}{x}
      \quad \text{变换方程}\\
      &\qquad
      \frac{\partial^2 z}{\partial x^2}
      -
      2\frac{\partial^2 z}{\partial x\partial y}
      +
      \frac{\partial^2 z}{\partial y^2}=0.
      \end{aligned}
      \]

\begin{hintbox}
这类题先写依赖图：\(z=e^{w+x+y}\) 或 \(z=xw\)，再把 \(\frac{\partial}{\partial x},\frac{\partial}{\partial y}\) 换成关于 \(u,v\) 的算子。
\end{hintbox}
\end{exercisebox}

\begin{exercisebox}
\textbf{习题 13\badge{含参数隐函数}}

设 \(y=f(x,t)\)，而 \(t\) 是由方程

\[
        F(x,y,t)=0
      \]

所确定的 \(x,y\) 的隐函数，其中 \(f\) 和 \(F\) 都具有连续偏导数。证明

\[
        \frac{dy}{dx}
        =
        \frac{
        \frac{\partial f}{\partial x}\frac{\partial F}{\partial t}
        -
        \frac{\partial f}{\partial t}\frac{\partial F}{\partial x}
        }{
        \frac{\partial F}{\partial t}
        -
        \frac{\partial f}{\partial t}\frac{\partial F}{\partial y}
        }.
      \]

\begin{hintbox}
这里 \(t=t(x,y)\)，而 \(y=f(x,t(x,y))\)。先对 \(y=f(x,t)\) 和 \(F(x,y,t)=0\) 同时求导，再消去 \(dt/dx\)。
\end{hintbox}
\end{exercisebox}

\begin{exercisebox}
\textbf{习题 14\badge{存在一阶积分}}

设二元函数 \(f(x,y):R^2\to R\) 具有连续偏导数。证明：存在一对一的连续向量值函数 \(G(t):R\to R^2\)，使得

\[
        f\circ G=\text{常数}.
      \]

\begin{hintbox}
从等值线角度理解：希望沿一条曲线走时 \(f\) 不变。可从常微分方程或梯度正交方向构造曲线。
\end{hintbox}
\end{exercisebox}

来源：陈纪修、于崇华、金路《数学分析》第三版下册，第十二章 §4，教材第 151-167 页。下一节：第十二章 §5：偏导数在几何中的应用。

上一节：第十二章 §3：中值定理和 Taylor 公式。复盘：错题集与好题集。路线图：教材路线图。
